\section{From Certified Relations to Positions}
\label{sec:portfolio}

\subsection{The edge representation}

Every active relation can be interpreted as a pair trade: long the predicted winner and short the predicted loser. Aggregate these edges through
\begin{equation}
\widetilde w_{i,t}(q)
=
\frac{1}{W_t}
\sum_{j\neq i}
\omega_{ij,t}
\left[
A_{ij,t}(q)-A_{ji,t}(q)
\right].
\label{eq:edge-portfolio}
\end{equation}
The portfolio is dollar neutral because the contribution of every edge sums to zero across its two endpoints. Its gross exposure declines endogenously when the investor abstains. This variable-notional portfolio measures the direct economic consequence of refusing uncertain distinctions.

For an equal-gross comparison, define
\[
w_{i,t}^{\mathrm{EG}}(q)
=
G
\frac{\widetilde w_{i,t}(q)}{
\sum_j|\widetilde w_{j,t}(q)|
}
\]
when the denominator is positive, where $G$ is the gross-leverage target. The variable-notional and equal-gross versions answer different questions. The first includes the decision to hold cash; the second isolates whether the surviving relations contain better information per unit of exposure.

The fixed denominator in \eqref{eq:edge-portfolio} gives an exact decomposition before rescaling:
\begin{equation}
\widetilde{\bm w}_t(0)
=
\widetilde{\bm w}_t(\widehat q_t)
+
\left[
\widetilde{\bm w}_t(0)
-
\widetilde{\bm w}_t(\widehat q_t)
\right].
\label{eq:exact-decomposition}
\end{equation}
The first component is certified and the second is the portfolio generated by removed relations. Equation \eqref{eq:exact-decomposition} permits a clean attribution of returns and changes in positions. Trading-cost attribution additionally accounts for the fact that costs are nonlinear in changes of the combined portfolio; we therefore report stand-alone costs, marginal costs, and a Shapley decomposition as robustness.

The paper also reports conventional top-minus-bottom decile portfolios. Those strategies are familiar benchmarks, but they do not give an exact relation-level decomposition. The edge portfolio is the primary economic experiment because it maps the certified object directly into capital.

\subsection{A no-trade interpretation}

Consider a single predicted relative-value opportunity. Let
\[
d_{ij,t}
=
\widehat\mu_{i,t}-\widehat\mu_{j,t}
\geq0,
\qquad
s_{ij,t}
=
\widehat s_{i,t}+\widehat s_{j,t}.
\]
The investor considers a nonnegative position $x$ that is long $i$ and short $j$. The unknown expected spread belongs to the ambiguity interval
\[
[d_{ij,t}-q s_{ij,t},
  d_{ij,t}+q s_{ij,t}].
\]
Let $c_{ij,t}\geq0$ be the proportional round-trip cost per unit of the pair position, and let $\gamma\nu_{ij,t}>0$ be the quadratic risk penalty.

\begin{proposition}[Certified ranking and the no-trade region]
\label{prop:no-trade}
The robust position
\begin{equation}
x_{ij,t}^{\star}(q)
=
\argmax_{x\geq0}
\min_{m\in[d_{ij,t}-q s_{ij,t},\,d_{ij,t}+q s_{ij,t}]}
\left\{xm-c_{ij,t}x-\frac{\gamma\nu_{ij,t}}{2}x^2\right\}
\label{eq:robust-pair-choice}
\end{equation}
is
\begin{equation}
x_{ij,t}^{\star}(q)
=
\frac{
[d_{ij,t}-q s_{ij,t}-c_{ij,t}]_+
}{\gamma\nu_{ij,t}}.
\label{eq:no-trade-solution}
\end{equation}
Consequently, the investor does not trade the predicted ranking relation whenever
\begin{equation}
d_{ij,t}
\leq
q s_{ij,t}+c_{ij,t}.
\label{eq:no-trade-region}
\end{equation}
\end{proposition}

Without costs, condition \eqref{eq:no-trade-region} contains the score-band overlap rule. Trading costs expand the no-trade region. The proposition is intentionally simple: it does not claim that every institutional portfolio solves independent pair problems. Its role is to show that abstention has a standard economic interpretation. Weak forecast differences should not automatically become positions when their worst-case spread is nonpositive net of implementation costs.

\subsection{Costs and capacity}

Let $\Delta w_{i,t}=w_{i,t}-w_{i,t-1}^{+}$ denote the trade after accounting for passive weight drift. We report a transparent linear benchmark
\[
\mathrm{TC}_{t}^{\mathrm{lin}}
=
c\sum_i|\Delta w_{i,t}|
\]
and a capacity-aware specification
\begin{equation}
\mathrm{TC}_{t}^{\mathrm{cap}}
=
\sum_i
\left[
\frac{\mathrm{spread}_{i,t}}{2}|\Delta w_{i,t}|
+
\lambda\sigma_{i,t}^{d}
\sqrt{
\frac{\mathrm{AUM}|\Delta w_{i,t}|}
{\mathrm{ADV}_{i,t}}
}
|\Delta w_{i,t}|
\right].
\label{eq:capacity-cost}
\end{equation}
The inputs use information available at portfolio formation. We trace results over AUM and risk rather than reporting one cost assumption. This follows the economic principle in \citet{JensenKellyMalamudPedersen2026}: a strategy should be judged on the attainable net-return frontier, not on gross alpha alone.
